\documentclass[11pt,a4paper]{article}

\usepackage[T1]{fontenc}
\usepackage[utf8]{inputenc}
\usepackage{lmodern}
\usepackage[margin=27mm]{geometry}
\usepackage{amsmath,amssymb,mathtools}
\usepackage[hidelinks]{hyperref}
\usepackage{microtype}

\newcommand{\R}{\mathbb{R}}
\newcommand{\ind}{\mathbf{1}}
\newcommand{\LN}{\operatorname{LN}}
\newcommand{\GN}{\operatorname{GN}}
\newcommand{\MHA}{\operatorname{MHA}}
\newcommand{\FFN}{\operatorname{FFN}}
\newcommand{\softmax}{\operatorname{softmax}}
\newcommand{\silu}{\operatorname{SiLU}}
\newcommand{\KL}{\operatorname{KL}}
\newcommand{\JS}{\operatorname{JS}}
\newcommand{\sg}{\operatorname{stopgrad}}
\newcommand{\clip}{\operatorname{clip}}
\newcommand{\BOS}{\mathrm{BOS}}
\newcommand{\EOS}{\mathrm{EOS}}
\newcommand{\eps}{\varepsilon}
\newcommand{\norm}[1]{\left\lVert #1\right\rVert}
\newcommand{\pos}[1]{\left[#1\right]_{+}}
\DeclareMathOperator{\concat}{concat}
\DeclareMathOperator{\Pool}{Pool}
\DeclareMathOperator{\Conv}{Conv}

\numberwithin{equation}{section}
\setlength{\parindent}{0pt}
\setlength{\parskip}{0.45em}
\setlength{\emergencystretch}{2em}
\allowdisplaybreaks[1]

\title{Causal Video-to-Text Learning\\
\large Mathematical Formulation and Training Procedure}
\author{}
\date{}

\begin{document}
\maketitle
\vspace{-2em}

\begin{abstract}
This document formulates the causal video-to-text model and two-stage
learning procedure defined in \texttt{train.py}. It describes aligned
supervision, causal visual and textual representations, multimodal
prediction, projection and attention regularization, and confidence targets
derived from a frozen predictor. The presentation is restricted to
mathematics and algorithmic operations.
\end{abstract}

The companion architecture document~\cite{architecture} follows the inputs
through every model block, with a flow diagram, tensor shapes, fusion
variants, and the cached streaming execution order. This document retains
the equations needed to define the training objectives; the architecture
document provides the detailed forward-pass reference.

\section{Sequences, alignments, and causal information}
\label{sec:alignment}

For one example, let $I_0,\ldots,I_{F-1}$ be video frames and
$y_1,\ldots,y_N$ the transcript tokens in a vocabulary $\mathcal V$ of
size $V>1$. Frame indices start at zero. Token $i$ has an inclusive
alignment window
\begin{equation}
\mathcal W_i=[a_i,e_i]\cap\mathbb Z,
\qquad 0\leq a_i\leq e_i<F,
\qquad e_i<a_{i+1}.
\end{equation}
Gaps between windows are permitted. Padding positions are excluded from
supervision and from attention keys.

The teacher sequence is $(\BOS,y_1,\ldots,y_N)$. Its causal text states are
\begin{equation}
z_0=E_t(\BOS),\qquad
z_j=E_t(\BOS,y_1,\ldots,y_j)_j,\quad 1\leq j\leq N.
\end{equation}
Thus the query for target $y_i$ is $q_i=z_{i-1}$: the target itself is
absent from its query. Token-query prediction also uses
$y_{N+1}=\EOS$, $q_{N+1}=z_N$, and the singleton window
$\mathcal W_{N+1}=\{F-1\}$. This additional target belongs only to the
true end of a video. Frame prediction and confidence supervision use
real-token windows only.

There are two distinct cross-attention masks. A token query may attend
to video up to the end of its own window:
\begin{equation}
M^{t\to v}_{if}=\ind\{0\leq f\leq e_i\}.
\label{eq:t2v}
\end{equation}
This mask allows all earlier frames, including frames before $a_i$.
Window-local attention is encouraged separately by a loss.

The video-to-text mask uses the time at which each text state becomes
available:
\begin{equation}
\alpha_0=0,\qquad \alpha_j=e_j+1\quad(j\geq1),
\qquad
M^{v\to t}_{fj}=\ind\{\alpha_j\leq f\}.
\label{eq:v2t}
\end{equation}
In particular, a token finishing at frame $e_j$ can influence this
branch only from frame $e_j+1$. Both masks are intersected with the
valid query and key positions for each example. Self-attention in
either encoder uses $M^{\mathrm{causal}}_{rs}=\ind\{s\leq r\}$, with
invalid keys removed.

\subsection{Alignment under tokenization}

If alignments are supplied for source characters, a token covering the
half-open character span $[c_i,d_i)$ receives the window
\begin{equation}
[a_i,e_i]=[a^{\mathrm{char}}_{c_i},
            e^{\mathrm{char}}_{d_i-1}].
\end{equation}
If $k$ consecutive tokens share exactly the same character span, its
$L$ frames are partitioned into $k$ consecutive nonempty intervals.
For part $r=0,\ldots,k-1$, the interval length is
\begin{equation}
L_r=\left\lfloor\frac{L}{k}\right\rfloor
       +\ind\{r<L\bmod k\},\qquad L\geq k.
\end{equation}
Token-aligned windows can instead be supplied directly. The resulting
token windows must remain ordered and non-overlapping.

\section{Causal representation learning}

For the block order and intermediate tensor shapes, see the video and text
path sections of the architecture document~\cite{architecture}.

All learned affine maps below have independent parameters unless sharing
is explicitly stated. Layer normalization $\LN$ acts on the feature
coordinates of one position. Group normalization $\GN$ in the visual
front end acts within a single frame, across spatial positions and
channel groups, and never aggregates across time.

Let $\sigma(x)=(1+e^{-x})^{-1}$ and $\silu(x)=x\sigma(x)$.
Write $\mathcal D_p$ for inverted dropout: during training,
$\mathcal D_p(x)=m\odot x/(1-p)$ with independent
$m_k\sim\operatorname{Bernoulli}(1-p)$; during deterministic
prediction it is the identity. Each occurrence draws its own mask.

\subsection{Masked multi-head attention}

For query vectors $u_r$ and key/value vectors $v_s$, head $h$ forms
\begin{equation}
Q_r^{h}=W_Q^{h}u_r+b_Q^{h},\qquad
K_s^{h}=W_K^{h}v_s+b_K^{h},\qquad
V_s^{h}=W_V^{h}v_s+b_V^{h}.
\end{equation}
For head width $d_h$, the pre-dropout probabilities are
\begin{equation}
A^{h}_{rs}=
\begin{cases}
\displaystyle
\frac{\exp(\langle Q_r^{h},K_s^{h}\rangle/\sqrt{d_h})}
 {\sum_{u:M_{ru}=1}\exp(\langle Q_r^{h},K_u^{h}\rangle/\sqrt{d_h})},
 & M_{rs}=1,\\[1.2em]
0, & M_{rs}=0.
\end{cases}
\label{eq:attention}
\end{equation}
A row with no admissible keys has all its probabilities set to zero.
With $H$ heads,
\begin{equation}
\MHA(U,V;M)_r
=W_O\concat_{h=1}^{H}\left(
 \sum_s\mathcal D_p(A^h)_{rs}V_s^h\right)+b_O.
\end{equation}
The attention regularizers below use $A^h$ before dropout.

The positionwise feed-forward transformation is
\begin{equation}
\FFN(x)=\mathcal D_p\!\left[
 W_2\mathcal D_p\!\left(\silu(W_1\LN(x)+b_1)\right)+b_2
\right].
\label{eq:ffn}
\end{equation}

\subsection{Visual front end}

A mouth crop $\mathcal C_{I_0}$ is fixed using the first video frame.
After resizing, the RGB intensities are normalized as
\begin{equation}
x_f=\frac{\mathcal C_{I_0}(I_f)}{127.5}-1.
\end{equation}
The same crop is retained throughout the video, so its selection does
not depend on future frames.

The first visual transformation is a causal spatiotemporal convolution:
\begin{equation}
u_f=\silu\!\left(\GN\!\left[
   \sum_{r=0}^{4}K_r*_{\mathrm{sp}}x_{f-r}
\right]\right),\qquad x_f=0\text{ for }f<0.
\end{equation}
Here $*_{\mathrm{sp}}$ is a spatial convolution. The spatial kernel is
$5\times5$ with stride two and centered padding; time is not
subsampled and padding is entirely in the past.

Each frame then passes through three spatial residual blocks. A block is
\begin{equation}
\mathcal R(u)=\silu\!\left(
 \GN_2\!\left[\Conv_2\bigl(\silu(\GN_1(\Conv_1(u)))\bigr)\right]
 +\mathcal S(u)\right),
\end{equation}
where the main convolutions have $3\times3$ kernels and $\mathcal S$
is the matching identity or $1\times1$ projection. The three blocks
have output widths $64$, $96$, and $128$, with spatial downsampling by
two in each block. Spatial averaging and an affine projection give
\begin{equation}
r_f=W_r\Pool\bigl(\mathcal R_3(\mathcal R_2(\mathcal R_1(u_f)))\bigr)+b_r
\in\R^{d_v}.
\end{equation}

\subsection{Positions and temporal visual layers}

For feature width $d$, let $h_d=\lceil d/2\rceil$ and
\begin{equation}
\omega_k=10000^{-k/\max(1,h_d-1)},\qquad
\mathrm{PE}_d(f)_{2k}=\sin(f\omega_k),\qquad
\mathrm{PE}_d(f)_{2k+1}=\cos(f\omega_k),
\label{eq:positions}
\end{equation}
retaining only coordinates smaller than $d$. The visual sequence begins
at $X^{(0)}_f=r_f+\mathrm{PE}_{d_v}(f)$.

One temporal visual layer applies
\begin{align}
U&=X+\tfrac12\FFN(X),\\
R&=U+\mathcal D_p\!\left[
 \MHA(\LN(U),\LN(U);M^{\mathrm{causal}})\right],\\
X^{+}&=R+\mathcal K(R).
\label{eq:visual-layer}
\end{align}
The causal convolutional transformation $\mathcal K$ is defined by
\begin{align}
[A_f,B_f]&=W_{\mathrm{in}}\LN(R_f)+b_{\mathrm{in}},\\
G_f&=A_f\odot\sigma(B_f),\\
\mathcal K(R)_f
 &=\mathcal D_p\!\left[
 W_{\mathrm{out}}\silu\!\left(
 b_{\mathrm{dw}}+\sum_{r=0}^{k_c-1}d_r\odot G_{f-r}\right)
 +b_{\mathrm{out}}\right].
\end{align}
The vectors $d_r$ specify a depthwise temporal convolution and
$G_f=0$ for negative indices. After $L_v$ such layers,
\begin{equation}
v_f=\LN(X_f^{(L_v)})\in\R^{d_v}.
\end{equation}
There is one half-weighted feed-forward residual in each visual layer,
followed by attention and convolution. Every $v_f$ depends only on
$x_0,\ldots,x_f$.

\subsection{Textual representations}

For a learned text encoder, set $t_0=\BOS$ and $t_j=y_j$ for $j\geq1$,
and initialize
\begin{equation}
Z_j^{(0)}=\operatorname{Embed}(t_j)+\mathrm{PE}_{d_t}(j).
\end{equation}
Each text layer is
\begin{align}
U&=Z+\mathcal D_p\!\left[
 \MHA(\LN(Z),\LN(Z);M^{\mathrm{causal}})\right],\\
Z^{+}&=U+\FFN(U).
\end{align}
After $L_t$ layers, $z_j=\LN(Z_j^{(L_t)})\in\R^{d_t}$.

Alternatively, $z_j$ is the final hidden state of a pretrained causal
language model on $(\BOS,y_1,\ldots,y_j)$. This alternative retains
that model's own internal architecture and positional representation;
its vocabulary prediction head is not used. Its parameters may be
fixed or optimized during the first training stage. A fixed text
encoder is evaluated deterministically.

\section{Multimodal prediction}

The two modalities are mapped to a common width $d$ by bias-free
projections (or the optional video identity mapping described below):
\begin{equation}
\bar v_f=P_vv_f,\qquad \bar z_j=P_tz_j,
\qquad \bar q_i=\bar z_{i-1}.
\label{eq:projections}
\end{equation}
By default, both projections are learned and $d=d_{\mathrm{fusion}}$.
With \texttt{--no-vproj}, $P_v$ is the identity, $d=d_v$, and $P_t$
maps text to $d_v$; all subsequent fusion operations use this width.
The earlier spatial encoder projection remains present. The architecture
document~\cite{architecture} details these connections and the two
prediction routes below.

\subsection{Token-query variants}

Text queries read the visual prefix through
\begin{equation}
c_i=\MHA(\bar Q,\bar V;M^{t\to v})_i.
\end{equation}
Let $\rho$ select the residual contribution of the text query:
\begin{equation}
\rho=
\begin{cases}
1,&\text{residual fusion},\\
0,&\text{cross-only fusion},\\
w_t,&\text{weighted fusion}.
\end{cases}
\end{equation}
The scalar $w_t$ is fixed, finite, and not learned. The shared
token-prediction transformation is
\begin{align}
h_i&=\LN_t(\rho\bar q_i+c_i),\\
\widehat h_i&=\LN_{t,\mathrm{out}}(h_i+\FFN_t(h_i)),\\
\ell_i&=W_y\widehat h_i+b_y,\qquad
p_i=\softmax(\ell_i).
\label{eq:token-head}
\end{align}
Consequently, the supervised token distribution conditions on the
transcript prefix and visual evidence through the aligned endpoint:
\begin{equation}
p_i(y)=p_\theta(y\mid y_{<i},x_{0:e_i}).
\end{equation}

\subsection{Frame prediction variant}

The frame variant has no text-to-video attention. Instead, each frame
queries the text states already available at that frame:
\begin{equation}
d_f=\MHA(\bar V,\bar Z;M^{v\to t})_f.
\label{eq:video-cross}
\end{equation}
It uses the token head on a visual residual:
\begin{align}
h_f^{\mathrm{fr}}&=\LN_t(\bar v_f+w_vd_f),\\
\widehat h_f^{\mathrm{fr}}
 &=\LN_{t,\mathrm{out}}\bigl(h_f^{\mathrm{fr}}+\FFN_t(h_f^{\mathrm{fr}})\bigr),\\
\ell_f^{\mathrm{fr}}&=W_y\widehat h_f^{\mathrm{fr}}+b_y,
\qquad p_f^{\mathrm{fr}}=\softmax(\ell_f^{\mathrm{fr}}).
\end{align}
Here $w_v$ is the fixed, finite \texttt{--vcross-weight} scalar, whose
default is one. At $w_v=0$, frame-token logits depend only on video.
This multiplier applies only to frame-token fusion; the confidence
branch below always uses $\bar v_f+d_f$.
For $f\in\mathcal W_i$, the current frame label is $y_i$.
Optional window phasing also supervises $y_{i-1}$ near the start of
the window, as defined in Section~\ref{sec:window-phasing}.
Frames in gaps have no supervised label. No frame is assigned an
$\EOS$ target.

\subsection{Confidence branch}

All variants produce a scalar confidence logit at each frame. With
$d_f$ defined by \eqref{eq:video-cross}, this branch applies
\begin{align}
g_f&=\LN_v(\bar v_f+d_f),\\
s_f^{v}&=\MHA(\LN_{v,a}(G),\LN_{v,a}(G);M^{\mathrm{causal}})_f,\\
g'_f&=\LN_{v,\mathrm{self}}(g_f+s_f^{v}),\\
\widehat g_f&=\LN_{v,\mathrm{out}}(g'_f+\FFN_v(g'_f)),\\
\kappa_f&=w_c^{\mathsf T}\widehat g_f+b_c.
\label{eq:confidence-head}
\end{align}
Here $G$ is the sequence of $g_f$ states. For temperature $\tau_c>0$,
the predicted confidence is
\begin{equation}
C_f=\sigma(\kappa_f/\tau_c).
\end{equation}
This prediction uses only the visual prefix and causally available text.

\section{Stage I: supervised prediction and regularization}

Batch index $b$ is restored in this section. Every average excludes
padding. Here $N_b$ counts the real-token windows assigned to the batch
element; for an unsplit video these are all its real-token windows.
A reduction over an empty eligible set contributes zero.

\subsection{Prediction loss}

Let $\mathcal I$ be the set of valid token queries in a batch, including
$\EOS$ queries at true video ends. Token-query prediction minimizes
\begin{equation}
\mathcal L_{\mathrm{pred}}
=-\frac{1}{|\mathcal I|}
 \sum_{(b,i)\in\mathcal I}\log p_{bi}(y_{bi}).
\label{eq:token-ce}
\end{equation}
For frame prediction, let $\mathcal J$ contain all batch/frame pairs
lying in a real-token window, including every nonempty intersection
with a video section (Section~\ref{sec:sections}). Let
$y_{b,f}^{\mathrm{fr}}$ be the unique token whose window contains
$f$. With window phasing disabled, the loss becomes
\begin{equation}
\mathcal L_{\mathrm{pred}}
=-\frac{1}{|\mathcal J|}
 \sum_{(b,f)\in\mathcal J}\log p_{bf}^{\mathrm{fr}}(y_{b,f}^{\mathrm{fr}}).
\label{eq:frame-ce}
\end{equation}
Thus a longer real-token window contributes more supervised positions
in the frame variant.

\subsection{Relative window phasing}
\label{sec:window-phasing}

Window phasing allows the previous token to remain a supervision target
over an initial portion of the current token's window. It applies only
to frame prediction. Let $\varphi\in[0,1]$ be the relative phasing
fraction, configured by \texttt{--window-phasing}; its default is zero.
Using the original, unclipped alignment window, define
\begin{equation}
L_{bi}=e_{bi}-a_{bi}+1,\qquad
k_{bi}=\lfloor\varphi L_{bi}\rfloor,\qquad
d_{bi}=k_{bi}-1.
\label{eq:phase-length}
\end{equation}
Here $k_{bi}$ is the number of phased frames and $d_{bi}$ is the last
included zero-based offset, not a frame count. The phased subset is
\begin{equation}
\mathcal P_{bi}=
\begin{cases}
\{a_{bi}+t: t\in\mathbb Z,\ 0\leq t\leq d_{bi}\},&i>1,\ k_{bi}>0,\\
\varnothing,&\text{otherwise}.
\end{cases}
\label{eq:phase-subset}
\end{equation}
The index $i$ refers to the original token sequence, so the first token
has no previous-token target; $\BOS$ is never added as a phased label.
For example, $\varphi=0.1$ and $L_{bi}=20$ phase offsets $0$ and $1$.
At this fraction a window shorter than ten frames has no phased frames.

For $f\in\mathcal W_{bi}$, the per-frame cross entropy is
\begin{equation}
\mathcal C_{bif}^{(\varphi)}=
\begin{cases}
-\tfrac12\log p_{bf}^{\mathrm{fr}}(y_{bi})
-\tfrac12\log p_{bf}^{\mathrm{fr}}(y_{b,i-1}),&f\in\mathcal P_{bi},\\
-\log p_{bf}^{\mathrm{fr}}(y_{bi}),&f\notin\mathcal P_{bi}.
\end{cases}
\label{eq:phased-frame-ce}
\end{equation}
This is cross entropy against a target assigning equal probability to
the current and previous tokens on phased frames. If those token IDs
are equal, their weights combine to one and the loss is ordinary
single-target cross entropy.

Let $i(b,f)$ identify the original token window containing a valid
frame. The first-stage prediction objective becomes
\begin{equation}
\mathcal L_{\mathrm{pred}}
=\frac{1}{|\mathcal J|}\sum_{(b,f)\in\mathcal J}
 \mathcal C_{b,i(b,f),f}^{(\varphi)},\qquad
|\mathcal J|=\sum_b|\{f:(b,f)\in\mathcal J\}|.
\label{eq:phased-batch-ce}
\end{equation}
Each valid frame counts once in the batch mean, including a phased
frame with two targets. Gaps and padding contribute neither loss nor
denominator mass. At $\varphi=0$, this reduces to
\eqref{eq:frame-ce}; at $\varphi=1$, every frame of each window after
the first uses both labels. Token-query objectives are unaffected.

Phasing changes supervision without shifting alignment boundaries or
teacher-token availability in \eqref{eq:v2t}. Stage~II retains the same
confidence windows and target construction, using the predictor learned
with the chosen phasing fraction. Evaluation continues to score the
current token $y_{bi}$ at each labeled frame.

\subsection{Projection geometry and norms}

For an eligible example, let the rows of $X\in\R^{L\times d_x}$ be
unprojected states and the rows of $H\in\R^{L\times d}$ their projected
counterparts. Define
\begin{equation}
\widetilde X_r=\frac{X_r}{\max(\norm{X_r}_2,\eps_n)},\qquad
\widetilde H_r=\frac{H_r}{\max(\norm{H_r}_2,\eps_n)},
\qquad \eps_n=10^{-6}.
\end{equation}
The geometry and norm penalties for that example are
\begin{align}
\mathcal G(X,H)
 &=\frac{1}{L^2}\norm{\widetilde X\widetilde X^{\mathsf T}
                     -\widetilde H\widetilde H^{\mathsf T}}_F^2,
\label{eq:geometry}\\
\mathcal N(X,H)
 &=\frac1L\sum_{r=1}^{L}
 \left(\frac{\norm{H_r}_2}{\norm{X_r}_2+\eps_n}-1\right)^2.
\label{eq:norm}
\end{align}
Equation~\eqref{eq:geometry} preserves pairwise cosine similarities;
Equation~\eqref{eq:norm} preserves individual feature magnitudes.

For video, use all valid $(v_f,\bar v_f)$ pairs. For text in the
token-query variants, use all supervised $(q_i,\bar q_i)$ pairs,
including a true-end $\EOS$ query. In the frame variant, use the full
valid teacher-state sequence $(z_j,\bar z_j)$. Only examples containing
at least one supervised prediction position are eligible. Averaging
\emph{per-example} penalties over these examples defines
$\mathcal L_{v,\mathrm{geom}}$, $\mathcal L_{t,\mathrm{geom}}$,
$\mathcal L_{v,\mathrm{norm}}$, and $\mathcal L_{t,\mathrm{norm}}$.
When \texttt{--no-vproj} is enabled, both video projection penalties
are skipped and contribute zero; text projection penalties are unchanged.

\subsection{Attention and cross-modal magnitude penalties}

These penalties apply only to token-query variants. Let
\begin{equation}
\overline A_{bif}=\frac1H\sum_{h=1}^{H}A^h_{bif},\qquad
\mu_{bi}=\sum_{f=0}^{F_b-1}f\,\overline A_{bif},
\end{equation}
where $A^h$ is the text-to-video attention from
\eqref{eq:attention}. For examples with at least two real-token
windows, define
\begin{equation}
\mathcal L_{\mathrm{mono}}
=\frac{1}{|\mathcal B_2|}\sum_{b\in\mathcal B_2}
 \frac{1}{N_b-1}\sum_{i=1}^{N_b-1}
 \pos{\mu_{bi}-\mu_{b,i+1}},
\quad \mathcal B_2=\{b:N_b\geq2\}.
\label{eq:monotonic}
\end{equation}
Here $[x]_+=\max(x,0)$. This discourages backward movement of the
attention mean; it does not require a strictly positive step.

Let $\mathcal W$ index all real-token windows in the batch. The
aligned-window penalty is
\begin{equation}
\mathcal L_{\mathrm{align}}
=-\frac{1}{|\mathcal W|}\sum_{(b,i)\in\mathcal W}
 \log\!\left(\max\!\left\{
     \sum_{f=a_{bi}}^{e_{bi}}\overline A_{bif},\,10^{-8}
 \right\}\right).
\label{eq:aligned}
\end{equation}
Both attention penalties exclude $\EOS$. Alignment averages over
windows, while monotonicity first averages adjacent pairs within each
eligible example.

For a prescribed margin $m\geq0$, the cross-modal magnitude penalty is
\begin{equation}
\mathcal L_{\mathrm{cross}}
=\frac1{|\mathcal I|}\sum_{(b,i)\in\mathcal I}
 \pos{m-\norm{c_{bi}}_2}^{\,2}.
\end{equation}
This penalty includes true-end $\EOS$ queries and discourages a
text-to-video contribution with norm below $m$.

\subsection{Combined first-stage objective}

The full objective is
\begin{align}
\mathcal L_{\mathrm I}={}&\mathcal L_{\mathrm{pred}}
 +\lambda_{v,g}\mathcal L_{v,\mathrm{geom}}
 +\lambda_{t,g}\mathcal L_{t,\mathrm{geom}}\notag\\
 &+\lambda_{v,n}\mathcal L_{v,\mathrm{norm}}
 +\lambda_{t,n}\mathcal L_{t,\mathrm{norm}}\notag\\
 &+\lambda_{\mathrm{mono}}\mathcal L_{\mathrm{mono}}
 +\lambda_{\mathrm{align}}\mathcal L_{\mathrm{align}}
 +\lambda_{\mathrm{cross}}\mathcal L_{\mathrm{cross}}.
\label{eq:stage-one}
\end{align}
The default objective has $\lambda_{v,n}=\lambda_{t,n}=0.01$ and all
other regularization coefficients equal to zero. The frame variant
requires $\lambda_{\mathrm{mono}}=\lambda_{\mathrm{align}}
=\lambda_{\mathrm{cross}}=0$, since it has no text-to-video branch.
Confidence-only transformations do not enter this objective.
When frame window phasing is enabled, $\mathcal L_{\mathrm{pred}}$ in
\eqref{eq:stage-one} is given by \eqref{eq:phased-batch-ce}.

\section{Stage II: confidence from a frozen predictor}

\subsection{Window-relative causal predictions}

After Stage~I, freeze every parameter affecting token or frame logits.
For a real-token window $i$, define
\begin{equation}
L_i=e_i-a_i+1,\qquad b_i=L_i-1,\qquad
f(i,t)=a_i+t,\quad 0\leq t\leq b_i.
\end{equation}
The symbol $b_i$ is the final relative timestep, not a batch index.

For token-query variants, hold $q_i$ fixed and reevaluate its prediction
at every frame in the window, replacing \eqref{eq:t2v} by
\begin{equation}
M^{t\to v,(i,t)}_{f}=\ind\{0\leq f\leq a_i+t\}.
\end{equation}
Denote the resulting frozen logits by $\ell^*_{i,t}$. For the frame
variant, $\ell^*_{i,t}$ is the frozen frame logit at $f=a_i+t$,
using the text-availability mask \eqref{eq:v2t}. These distributions
use teacher-forced text history with commits at aligned window ends.

With token temperature $\tau_y>0$, form
\begin{equation}
P_{i,t}=\softmax(\ell^*_{i,t}/\tau_y),\qquad
Y_{i,t,v}=\frac{\max(P_{i,t,v},\eps_y)}
 {\sum_{u=1}^{V}\max(P_{i,t,u},\eps_y)},
\quad \eps_y=10^{-8}.
\label{eq:target-dist}
\end{equation}
The coordinates $(i,t)$ range only over real-token windows; there
are no confidence targets for gaps, padding, or $\EOS$.

\subsection{Uncertainty and future-distribution instability}

All logarithms are natural. Normalized token entropy is
\begin{equation}
U_{i,t}=\frac{-\sum_{v=1}^{V}Y_{i,t,v}\log Y_{i,t,v}}{\log V}.
\label{eq:entropy}
\end{equation}
For every admissible future offset $u\in\{0,\ldots,b_i-t\}$, set
\begin{align}
M_{i,t,u}&=\tfrac12(Y_{i,t}+Y_{i,t+u}),\\
D_{i,t,u}&=\JS(Y_{i,t},Y_{i,t+u})\notag\\
 &=\tfrac12\KL(Y_{i,t}\Vert M_{i,t,u})
   +\tfrac12\KL(Y_{i,t+u}\Vert M_{i,t,u}),
\label{eq:js}
\end{align}
where $\KL(p\Vert q)=\sum_v p_v\log(p_v/q_v)$.
Then $0\leq D_{i,t,u}\leq\log2$ and $D_{i,t,0}=0$.

For decay $\gamma\geq0$, the future-offset weights are
\begin{equation}
w_{i,t,u}=\frac{e^{-\gamma u}}
 {\sum_{r=0}^{b_i-t}e^{-\gamma r}},\qquad
\sum_{u=0}^{b_i-t}w_{i,t,u}=1.
\end{equation}
The default $\gamma=0$ gives uniform weights, including the
current distribution at $u=0$. With $\beta>0$, the mathematical
instability aggregate is
\begin{equation}
Q_{i,t}=\frac1\beta\log\left(
 \sum_{u=0}^{b_i-t}w_{i,t,u}e^{\beta D_{i,t,u}}\right).
\label{eq:instability}
\end{equation}
For normalized weights, $0\leq Q_{i,t}\leq\log2$.
For fixed positive weights its limiting forms are
\begin{equation}
\lim_{\beta\to0^+}Q_{i,t}
 =\sum_u w_{i,t,u}D_{i,t,u},\qquad
\lim_{\beta\to\infty}Q_{i,t}=\max_u D_{i,t,u}.
\end{equation}

The evaluated aggregate uses a logarithmic weight floor:
\begin{equation}
Q^{\eps}_{i,t}=\frac1\beta\log\left(
 \sum_{u=0}^{b_i-t}\max(w_{i,t,u},10^{-8})
                         e^{\beta D_{i,t,u}}\right).
\label{eq:instability-floor}
\end{equation}
The floored weights are not renormalized. Therefore the target includes
clipping even though the unfloored normalized quantities are bounded.
For mixing coefficient $\lambda\in[0,1]$, define
\begin{equation}
S_{i,t}=\clip_{[0,1]}\!\left(
  \lambda U_{i,t}+(1-\lambda)\frac{Q^{\eps}_{i,t}}{\log2}
\right),\qquad
s_{i,t}=1-S_{i,t}.
\label{eq:soft-target}
\end{equation}
Here $\lambda$ is the entropy/instability mixture coefficient, distinct
from the Stage~I regularization coefficients. The defaults are
$\beta=10$, $\lambda=0.5$, $\gamma=0$, and $\tau_y=1$.

At the last timestep, only $u=0$ remains, so $Q^{\eps}_{i,b_i}=0$
and $s_{i,b_i}=1-\lambda U_{i,b_i}$. The target measures predictive
certainty and stability; it does not directly compare the predicted
token with the true token. Future distributions enter target
construction only. The confidence predictor itself remains causal.

\subsection{Confidence loss and parameter separation}

Pack each frame confidence into the same window-relative coordinates:
\begin{equation}
C_{bi,t}=\sigma(\kappa_{b,a_{bi}+t}/\tau_c),\qquad
\mathcal T=\{(b,i,t):i\text{ is real},\ 0\leq t\leq b_{bi}\}.
\end{equation}
With $\eps_c=10^{-7}$, the optimized loss is
\begin{align}
\mathcal L_{\mathrm{II}}
=-\frac1{|\mathcal T|}\sum_{(b,i,t)\in\mathcal T}
\bigl[&\sg(s_{bi,t})\log\max(C_{bi,t},\eps_c)\notag\\
 &+(1-\sg(s_{bi,t}))\log\max(1-C_{bi,t},\eps_c)\bigr].
\label{eq:confidence-loss}
\end{align}
The default confidence temperature is $\tau_c=1$. This is an average
over valid timesteps, so longer windows have greater total weight.

Let $\theta$ contain all parameters that affect token/frame predictions
and $\phi$ contain the parameters exclusive to confidence prediction.
During Stage~II,
\begin{equation}
\theta=\theta^*,\qquad
\min_{\phi}\mathcal L_{\mathrm{II}}(\phi;\theta^*),\qquad
\nabla_\phi s_{bi,t}=0.
\end{equation}
Both encoders, both modality projections, and the token-prediction
transformations stay fixed and deterministic. The video-side
normalizations, causal self-attention, feed-forward transformation,
and scalar confidence head are trainable. Video-to-text attention is
also trainable in token-query variants; in the frame variant it must
remain fixed because it also determines frame-token logits. Dropout
is active only in the trainable confidence transformations.

\section{Algorithmic training and sequence sections}

\subsection{Optimization procedure}

Stage~I repeatedly constructs the aligned causal predictions and
minimizes \eqref{eq:stage-one}. An imported language model remains
fixed unless its fine-tuning is selected. Once prediction training
finishes, Stage~II constructs \eqref{eq:target-dist}--\eqref{eq:soft-target}
with the frozen predictor and minimizes \eqref{eq:confidence-loss} over
the confidence parameters. A batch with no eligible supervised
positions is skipped for the corresponding stage. Updating the
prediction parameters changes the distributions defining the
confidence targets and therefore requires fitting confidence again.

Both stages use gradient norm clipping and AdamW, with separate
optimization states. For the active parameter vector $\vartheta$ and
loss $\mathcal L$, clipping at threshold $c>0$ gives
\begin{equation}
g_k=\nabla_{\vartheta}\mathcal L(\vartheta_{k-1}),\qquad
\widetilde g_k=g_k\min\!\left(1,\frac{c}{\norm{g_k}_2+10^{-6}}\right).
\end{equation}
With moments initialized at zero, the update is
\begin{align}
m_k&=\beta_1m_{k-1}+(1-\beta_1)\widetilde g_k,\\
r_k&=\beta_2r_{k-1}+(1-\beta_2)\widetilde g_k^{\odot2},\\
\widehat m_k&=m_k/(1-\beta_1^k),\qquad
\widehat r_k=r_k/(1-\beta_2^k),\\
\vartheta_k
 &=(1-\eta\omega)\vartheta_{k-1}
   -\eta\,\frac{\widehat m_k}{\sqrt{\widehat r_k}+\eps_{\mathrm{opt}}}.
\end{align}
Products, square roots, and division in the adaptive term are
coordinatewise. Here $\eta$ is the stage-specific learning rate and
$\omega$ the decoupled weight-decay coefficient. The moment parameters
are $\beta_1=0.9$, $\beta_2=0.999$, and $\eps_{\mathrm{opt}}=10^{-8}$.
Frozen parameters and parameters outside the active loss path receive
no update.

\subsection{Optional division into video sections}
\label{sec:sections}

If a video is divided into consecutive sections $[r,s)$, the previous
equations apply with local frame index $f'=f-r$ and the following
supervision rules. Token-query targets are assigned exactly once, to
the section containing their endpoint:
\begin{equation}
\mathcal I_{[r,s)}=\{i:r\leq e_i<s\},\qquad
[a_i',e_i']=[\max(a_i,r)-r,\ e_i-r].
\end{equation}
Let $K=|\{i:e_i<r\}|$ and $M=|\mathcal I_{[r,s)}|$.
The text encoder still receives
$(\BOS,y_1,\ldots,y_{K+M})$, and the selected prediction queries are
$z_K,\ldots,z_{K+M-1}$. If $s=F$, also include the $\EOS$ query
$z_{K+M}$ with local window $\{s-r-1\}$.
Text availability is translated, rather than reset:
\begin{equation}
\alpha'_0=-r,\qquad \alpha'_j=e_j+1-r.
\end{equation}
Previously committed text may consequently have negative availability
times and is visible from the start of the section.

Frame prediction uses every nonempty intersection
\begin{equation}
\mathcal W_i\cap[r,s),\qquad
[a_i',e_i']=[\max(a_i,r)-r,\ \min(e_i,s-1)-r],
\end{equation}
including a token that finishes in a later section. Confidence uses
the clipped assigned windows for token-query variants, and all
nonempty window intersections for the frame variant. Its future
offsets terminate at the end of that local window.

Window phasing is computed before clipping. In section-local coordinates,
the phased subset is
\begin{equation}
\mathcal P'_{i,[r,s)}
=\{f-r:f\in\mathcal P_i\cap[r,s)\}.
\label{eq:section-phasing}
\end{equation}
Thus $k_i=\lfloor\varphi(e_i-a_i+1)\rfloor$ always uses the full
original window length, and the previous label remains the original
$y_{i-1}$ even if its window lies in an earlier section. Phasing never
restarts at a section boundary. If a cut falls inside the phased subset,
only its remaining frames retain the additional target.

The full preceding text prefix is retained, but visual temporal
context and visual positions begin anew in each section. The mouth
crop remains the one selected at the original video's first frame.
Thus section-based learning has a shorter visual context than
full-video learning, and clipping may also shorten the horizon used
to construct confidence targets.

\subsection{Causal sequential realization}

For the step-by-step streaming calls and the caches retained by each
block, see the execution-order section of the architecture
document~\cite{architecture}.

Within the chosen visual context, frame $f$ produces $v_f$ from frames
through $f$ and combines it with text states committed before its
video-to-text attention step. A token query predicts from all visual
states observed so far. Appending a committed token creates a new
causal text state, which becomes available to subsequent video steps.
Convolutions retain only past inputs, and attention admits only the
keys allowed by the corresponding causal mask.

With deterministic transformations, these sequential operations
evaluate the same causal functions as the masked equations for the
same video context and text-commit history. Confidence target
construction uses the aligned teacher-forced history; a generated
history with different tokens or commit times defines different
conditioning information.

\begin{thebibliography}{9}
\bibitem{architecture}
\emph{Causal Video-to-Text Learning: Model Architecture and Forward-Pass
Sequence}. Companion document, \texttt{model\_architecture.tex}.
\href{model_architecture.pdf}{Compiled architecture document}.
\end{thebibliography}

\end{document}
